% !TEX root = ../main.tex
\chapter*{Conclusion and outlook}
\phantomsection
\addcontentsline{toc}{chapter}{Conclusion and outlook}
\markboth{Conclusion}{Conclusion}
\label{sec:conclusion}

This internship project investigated how multi-hop question answering can be
made more adaptive and evidence-aware. The work addressed the problem from two
complementary perspectives: first, whether the compositional difficulty of a
question can be identified before answer generation; and second, once
decomposition is used, how the resulting retrieval process can be evaluated and
improved.

The first part of the thesis showed that benchmark-defined compositional
difficulty is linearly decodable from the internal representations of several
language models. Across three models and three multi-hop benchmarks, the signal
emerged across network depth, remained above a label-permutation null, survived
partial control for selected surface features, exhibited ordinal structure, and
transferred more consistently across benchmarks after surface residualisation.
These results provide evidence that question representations contain a graded
component associated with compositional demand, which could ultimately provide
a useful signal for adaptive processing before the full cost of answering is
incurred.

Once decomposition is employed, the thesis focused on the evidence retrieved by
its subquestions. An evidence-aware evaluation framework was developed to
separate final-answer correctness from evidence availability and retrieval
depth. The framework also introduced an evidence-ownership estimator for
associating supporting passages with the subquestions expected to retrieve
them. On the reference-style control, 92 of 100 complete ownership masks were
recovered. On the 100-question MuSiQue development panel, generated
decomposition increased evidence availability and, with GPT-4o-mini, raised
exact match from 40\% to 52\%. The experiments also showed that improved
retrieval does not necessarily produce an equivalent answer-quality gain,
highlighting the importance of intermediate answers and downstream reasoning in
the complete execution.

The final part of the work investigated how these measurements could support
optimisation. Automated search over decomposition instructions demonstrated that
the retrieval-oriented score can be used as an optimisation objective, although
the selected prompt achieved only a modest development-set improvement and did
not improve exact match. More importantly, the candidate comparison showed that
the prompt favoured by retrieval depth was not the one favoured by final-answer
quality, illustrating that retrieval-oriented and end-to-end objectives are not
fully interchangeable. A learned passage allocator was also developed to
distribute a fixed retrieval budget across subquestions. Under the tight
eight-slot setting, unit-increment allocation increased complete-evidence
coverage by 6.3 percentage points on 2WikiMultiHopQA and 4.7 points on MuSiQue,
while a separate 700-question MuSiQue experiment produced a 2.1-point F1 gain
for the learned segment policy. The advantage diminished at larger budgets,
indicating that adaptive allocation is particularly relevant when context is
constrained.

Several limitations remain. The difficulty study relies on benchmark-defined
proxies and establishes decodability rather than a causal role of the identified
representation. Ownership estimation on naturally generated decompositions
requires broader independent validation. The prompt-search pilot uses a shared
development panel, while the passage-allocation experiments operate on recorded
retrieval traces rather than fully regenerated sequential executions. The
structured evidence-feedback interface developed in Chapter~5 has been
diagnostically validated, but its comparative effect on prompt optimisation
remains to be evaluated.

These limitations define the main directions for future work. A natural next
step is to connect the pre-answer difficulty signal with an adaptive routing
policy that considers not only predicted difficulty, but also the expected
benefit and computational cost of decomposition. Passage allocation should also
be extended from fixed-trace replay to online sequential execution, where each
budget decision can affect intermediate answers and subsequent queries.
Finally, the proposed independent evaluation can determine whether structured
evidence feedback improves prompt search compared with answer-only or raw
evidence feedback.

Overall, this thesis shows that adaptive multi-hop question answering can be
studied at several complementary levels: compositional demand can be detected
before answering, decomposition can be evaluated through the evidence it
exposes, and retrieval resources can be allocated according to the needs of
individual subquestions rather than uniformly. These results provide a
foundation for future RAG systems in which reasoning and retrieval behaviour
are increasingly adapted to the structure and evidence requirements of each
question.